\documentclass{article}
\usepackage{amsmath,amssymb,amsthm}
\usepackage{algorithm,algorithmic}
\usepackage{hyperref}
\usepackage{enumitem}
\usepackage{geometry}
\geometry{margin=1in}
\newtheorem{theorem}{Theorem}
\newtheorem{lemma}{Lemma}
\newtheorem{assumption}{Assumption}
\newcommand{\E}{\mathbb{E}}
\newcommand{\R}{\mathbb{R}}
\newcommand{\norm}[1]{\left\lVert#1\right\rVert}
\newcommand{\inner}[1]{\left\langle#1\right\rangle}
\begin{document}

\section*{Algorithm Name}
\textbf{SAGA for Non‑Convex Finite‑Sum Optimization (SAGA‑NC).}

\section*{Mathematical Formulation}
Consider the finite‑sum objective  
\[
f(x)\;=\;\frac1n\sum_{i=1}^{n}f_i(x),\qquad x\in\R^{d},
\]
where each component $f_i$ is \emph{L‑smooth} (i.e. $\|\nabla f_i(x)-\nabla f_i(y)\|\le L\|x-y\|$) but not necessarily convex.  

Maintain a \emph{gradient table} 
\[
\phi_i^k\in\R^{d},\qquad g_i^k:=\nabla f_i(\phi_i^k),\qquad i=1,\dots,n,
\]
initialized arbitrarily (commonly $\phi_i^0=x_0$).  
At iteration $k$ draw an index $i_k\sim\text{Uniform}\{1,\dots,n\}$ and construct the unbiased variance‑reduced estimator
\[
v_k\;:=\;\underbrace{\nabla f_{i_k}(x_k)-\nabla f_{i_k}(\phi_{i_k}^k)}_{\text{correction}}
\;+\;\underbrace{\frac1n\sum_{j=1}^{n}g_j^k}_{\text{average table gradient}}.
\tag{1}
\]
The iterate is updated by a constant stepsize $\eta>0$:
\[
x_{k+1}\;=\;x_k\;-\;\eta\,v_k. \tag{2}
\]
Finally the table entry for the sampled index is refreshed:
\[
\phi_{i_k}^{k+1}\;=\;x_k,\qquad 
g_{i_k}^{k+1}\;=\;\nabla f_{i_k}(x_k), \qquad 
\phi_{j}^{k+1}= \phi_{j}^{k},\; g_{j}^{k+1}=g_{j}^{k}\;\;(j\neq i_k). \tag{3}
\]

\section*{Pseudocode}
\begin{algorithm}[H]
\caption{SAGA for Non‑Convex Finite‑Sum}
\begin{algorithmic}[1]
\STATE \textbf{Input:} stepsize $\eta>0$, initial point $x_0$, number of iterations $T$.
\STATE Initialize $\phi_i^0\leftarrow x_0$ and $g_i^0\leftarrow \nabla f_i(x_0)$ for all $i$.
\STATE Compute $\bar g^0\leftarrow \frac1n\sum_{i=1}^{n}g_i^0$.
\FOR{$k=0,\dots,T-1$}
    \STATE Sample $i_k\sim\text{Uniform}\{1,\dots,n\}$.
    \STATE $v_k\leftarrow \nabla f_{i_k}(x_k)-g_{i_k}^k+\bar g^k$.
    \STATE $x_{k+1}\leftarrow x_k-\eta v_k$.
    \STATE Update the table entry:
    \[
    g_{i_k}^{k+1}\leftarrow \nabla f_{i_k}(x_k),\qquad 
    \phi_{i_k}^{k+1}\leftarrow x_k .
    \]
    \STATE Update the average gradient:
    \[
    \bar g^{k+1}\leftarrow \bar g^{k}+ \frac{1}{n}\bigl(g_{i_k}^{k+1}-g_{i_k}^{k}\bigr).
    \]
\ENDFOR
\STATE \textbf{Output:} $\{x_k\}_{k=0}^{T}$ (or a uniform random iterate).
\end{algorithmic}
\end{algorithm}

\section*{Dry Run Example}
We minimise the simple scalar function 
\[
f(w)=(w-3)^2,\qquad f_i\equiv f\;(n=1).
\]
Hence $\nabla f(w)=2(w-3)$.  Choose $w_0=0$, stepsize $\eta=0.1$, and run three iterations.

\begin{center}
\begin{tabular}{c|c|c|c}
Iteration $k$ & $w_k$ & $v_k$ (gradient estimator) & $f(w_k)$\\\hline
0 & $0$ & $2(0-3)=-6$ & $9$\\
1 & $w_1=0-0.1(-6)=0.6$ & $2(0.6-3)=-4.8$ & $5.76$\\
2 & $w_2=0.6-0.1(-4.8)=1.08$ & $2(1.08-3)=-3.84$ & $3.6864$\\
3 & $w_3=1.08-0.1(-3.84)=1.464$ & $2(1.464-3)=-3.072$ & $2.3689$
\end{tabular}
\end{center}
Because $n=1$, the table always stores the current point, so the estimator $v_k$ coincides with the true gradient.  The objective value decreases steadily, illustrating the mechanics of (1)–(3).

\newpage
\section*{Assumptions}
\begin{assumption}[Smoothness]\label{as:smooth}
Each component $f_i$ is $L$‑smooth:
\[
\forall x,y\in\R^{d},\qquad 
\|\nabla f_i(x)-\nabla f_i(y)\|\le L\|x-y\|.
\]
\end{assumption}
\begin{assumption}[Bounded Gradient at Optimum]\label{as:bounded}
Define a (possibly non‑unique) stationary point $x^{\star}$ satisfying $\nabla f(x^{\star})=0$.  
Assume the average of squared gradients at $x^{\star}$ is bounded:
\[
\sigma^{2}\;:=\;\frac1n\sum_{i=1}^{n}\|\nabla f_i(x^{\star})\|^{2}<\infty .
\]
\end{assumption}
\begin{assumption}[Stepsize]\label{as:stepsize}
The constant stepsize obeys 
\[
0<\eta\le\frac{1}{3L}.
\]
\end{assumption}

\section*{Main Convergence Theorem}
\begin{theorem}[Non‑Convex SAGA Convergence]\label{thm:main}
Under Assumptions~\ref{as:smooth}–\ref{as:stepsize}, let $\{x_k\}_{k\ge0}$ be the iterates generated by SAGA‑NC.  
Define the random output $\tilde x$ as a uniform draw from $\{x_0,\dots,x_{T-1}\}$.  
Then
\[
\E\!\left[\bigl\|\nabla f(\tilde x)\bigr\|^{2}\right]
\;\le\;
\underbrace{\frac{2\bigl(f(x_0)-f(x^{\star})\bigr)}{\eta T}}_{\text{optimization term}}
\;+\;
\underbrace{4L\eta\,\sigma^{2}}_{\text{variance term}}.
\tag{4}
\]
Consequently, choosing $\eta = \Theta\!\bigl(1/\sqrt{T}\bigr)$ yields the optimal $O(1/\sqrt{T})$ rate for the expected squared gradient norm.
\end{theorem}

\section*{Theoretical Properties}
\begin{itemize}[leftmargin=*]
\item \textbf{Stability.} The bound (4) is linear in $\eta$ for the variance term and inversely proportional to $\eta$ for the optimization term; the stepsize condition $\eta\le 1/(3L)$ guarantees that the descent lemma (Lemma~\ref{lem:descent}) holds, preventing divergence.
\item \textbf{Hyper‑parameter Sensitivity.} The optimal $\eta$ balances the two terms in (4).  Too small $\eta$ inflates the $1/(\eta T)$ term, whereas too large $\eta$ violates Assumption~\ref{as:stepsize} and destroys the contraction used in the proof.
\item \textbf{Comparison to Baselines.}
    \begin{enumerate}
        \item \emph{Stochastic Gradient Descent (SGD).} For non‑convex $L$‑smooth problems, SGD with stepsize $\eta_t=\Theta(1/\sqrt{t})$ achieves $\E\|\nabla f(\tilde x)\|^2=O(1/\sqrt{T})$ but with a larger constant due to variance $\sigma_{\text{SGD}}^2$.  SAGA reduces the variance term to $\sigma^{2}$, often much smaller because stored gradients track past iterates.
        \item \emph{SVRG.} SVRG attains the same $O(1/\sqrt{T})$ rate but requires an inner loop and full‑gradient recomputation.  SAGA is fully online and enjoys the same theoretical guarantee without epoch resets.
    \end{enumerate}
\end{itemize}

\section*{Discussion}
The theorem shows that, despite the lack of convexity, SAGA‑NC drives the expected gradient norm to zero at the same $1/\sqrt{T}$ speed as classical SGD, yet with a variance term that depends only on the gradients at a stationary point rather than on the full stochastic noise.  The proof hinges on the unbiasedness of the estimator (1) and a careful control of its second moment using the smoothness of each $f_i$.

\newpage
\section*{Preliminaries (Definitions and Lemmas)}
\begin{lemma}[Smoothness Inequality]\label{lem:smooth}
If $f_i$ is $L$‑smooth then for any $x,y$,
\[
f_i(y)\le f_i(x)+\langle\nabla f_i(x),y-x\rangle+\frac{L}{2}\|y-x\|^{2}.
\]
\end{lemma}
\begin{proof}
Standard consequence of $L$‑smoothness; see e.g. Nesterov (2004). \hfill $\square$
\end{proof}

\begin{lemma}[Unbiasedness of SAGA Estimator]\label{lem:unbiased}
For any $k$, conditional on the sigma‑algebra $\mathcal{F}_k:=\sigma\{x_0,\phi_i^0,\dots,\phi_i^k\}$,
\[
\E_{i_k}[v_k\mid\mathcal{F}_k]=\nabla f(x_k).
\]
\end{lemma}
\begin{proof}
\[
\E_{i_k}[v_k\mid\mathcal{F}_k]
=\frac1n\sum_{i=1}^{n}\bigl(\nabla f_i(x_k)-\nabla f_i(\phi_i^k)\bigr)
+\frac1n\sum_{j=1}^{n}\nabla f_j(\phi_j^k)
=\frac1n\sum_{i=1}^{n}\nabla f_i(x_k)=\nabla f(x_k).
\] \hfill $\square$
\end{proof}

\begin{lemma}[Second‑Moment Bound]\label{lem:second}
Define $D_k:=\frac1n\sum_{i=1}^{n}\|x_k-\phi_i^k\|^{2}$.  Then
\[
\E_{i_k}\!\bigl[\|v_k\|^{2}\mid\mathcal{F}_k\bigr]
\;\le\;
4L^{2}D_k\;+\;2\|\nabla f(x_k)\|^{2}.
\tag{5}
\]
\end{lemma}
\begin{proof}
\[
\begin{aligned}
\|v_k\|^{2}
&=\Bigl\|\underbrace{\nabla f_{i_k}(x_k)-\nabla f_{i_k}(\phi_{i_k}^k)}_{\Delta_{i_k}}
+\underbrace{\frac1n\sum_{j=1}^{n}\nabla f_j(\phi_j^k)}_{\bar g^k}\Bigr\|^{2}\\[4pt]
&\le 2\|\Delta_{i_k}\|^{2}+2\Bigl\|\bar g^k\Bigr\|^{2}
\qquad\text{(by } \|a+b\|^{2}\le2\|a\|^{2}+2\|b\|^{2}\text{)}. \tag{Step 1}
\end{aligned}
\]
Using $L$‑smoothness,
\[
\|\Delta_{i_k}\|
= \|\nabla f_{i_k}(x_k)-\nabla f_{i_k}(\phi_{i_k}^k)\|
\le L\|x_k-\phi_{i_k}^k\|,
\]
hence $\|\Delta_{i_k}\|^{2}\le L^{2}\|x_k-\phi_{i_k}^k\|^{2}$.  Taking expectation over $i_k$,
\[
\E_{i_k}[\|\Delta_{i_k}\|^{2}\mid\mathcal{F}_k]
= \frac1n\sum_{i=1}^{n}L^{2}\|x_k-\phi_i^k\|^{2}
= L^{2}D_k. \tag{Step 2}
\]
For the second term,
\[
\bar g^k = \frac1n\sum_{j=1}^{n}\nabla f_j(\phi_j^k)
= \nabla f(x_k)-\frac1n\sum_{j=1}^{n}\bigl(\nabla f_j(x_k)-\nabla f_j(\phi_j^k)\bigr).
\]
Thus,
\[
\|\bar g^k\|^{2}
\le 2\|\nabla f(x_k)\|^{2}
+2\Bigl\|\frac1n\sum_{j=1}^{n}\bigl(\nabla f_j(x_k)-\nabla f_j(\phi_j^k)\bigr)\Bigr\|^{2}.
\]
Using Jensen’s inequality,
\[
\Bigl\|\frac1n\sum_{j}\bigl(\nabla f_j(x_k)-\nabla f_j(\phi_j^k)\bigr)\Bigr\|^{2}
\le \frac1n\sum_{j}\|\nabla f_j(x_k)-\nabla f_j(\phi_j^k)\|^{2}
\le L^{2}D_k. \tag{Step 3}
\]
Combining the bounds from Steps 1–3,
\[
\E_{i_k}[\|v_k\|^{2}\mid\mathcal{F}_k]
\le 2L^{2}D_k+2\bigl(2\|\nabla f(x_k)\|^{2}+L^{2}D_k\bigr)
=4L^{2}D_k+2\|\nabla f(x_k)\|^{2}.
\] \hfill $\square$
\end{proof}

\begin{lemma}[Descent Lemma for SAGA]\label{lem:descent}
Let $x_{k+1}=x_k-\eta v_k$ with $0<\eta\le 1/(3L)$.  Then
\[
\E_{i_k}\!\bigl[f(x_{k+1})\mid\mathcal{F}_k\bigr]
\le f(x_k)-\frac{\eta}{2}\|\nabla f(x_k)\|^{2}
+\frac{2\eta L^{2}}{n}\sum_{i=1}^{n}\|x_k-\phi_i^k\|^{2}.
\tag{6}
\]
\end{lemma}
\begin{proof}
By Lemma~\ref{lem:smooth} applied to $f$ (the average of the $f_i$):
\[
f(x_{k+1})\le f(x_k)+\langle\nabla f(x_k),x_{k+1}-x_k\rangle+\frac{L}{2}\|x_{k+1}-x_k\|^{2}.
\]
Insert the update (2):
\[
f(x_{k+1})\le f(x_k)-\eta\langle\nabla f(x_k),v_k\rangle+\frac{L\eta^{2}}{2}\|v_k\|^{2}. \tag{Step 4}
\]
Take expectation conditioned on $\mathcal{F}_k$ and use Lemma~\ref{lem:unbiased}:
\[
\E_{i_k}[\langle\nabla f(x_k),v_k\rangle\mid\mathcal{F}_k]
= \|\nabla f(x_k)\|^{2}. \tag{Step 5}
\]
Apply Lemma~\ref{lem:second} to the quadratic term:
\[
\E_{i_k}[\|v_k\|^{2}\mid\mathcal{F}_k]
\le 4L^{2}D_k+2\|\nabla f(x_k)\|^{2}. \tag{Step 6}
\]
Plugging Steps 5–6 into Step 4 gives
\[
\E_{i_k}[f(x_{k+1})\mid\mathcal{F}_k]
\le f(x_k)-\eta\|\nabla f(x_k)\|^{2}
+\frac{L\eta^{2}}{2}\bigl(4L^{2}D_k+2\|\nabla f(x_k)\|^{2}\bigr).
\]
Rearrange:
\[
= f(x_k)-\Bigl(\eta-\eta^{2}L\Bigr)\|\nabla f(x_k)\|^{2}
+2L^{3}\eta^{2}D_k.
\]
Since $\eta\le 1/(3L)$, we have $\eta-\eta^{2}L\ge \eta/2$.  Moreover $2L^{3}\eta^{2}=2\eta L^{2}\cdot(L\eta)\le 2\eta L^{2}$ because $L\eta\le 1/3<1$.  Hence
\[
\E_{i_k}[f(x_{k+1})\mid\mathcal{F}_k]
\le f(x_k)-\frac{\eta}{2}\|\nabla f(x_k)\|^{2}
+2\eta L^{2}D_k,
\]
which is exactly (6). \hfill $\square$
\end{proof}

\section*{Main Proof (Step‑by‑Step Derivation)}
We now prove Theorem~\ref{thm:main}.  All expectations are taken with respect to the whole random trajectory $\{i_0,\dots,i_{T-1}\}$ unless otherwise noted.

\begin{enumerate}[label={[\textbf{Step \arabic*}]}]
\item \textbf{Summation of Descent Inequality.}  
Take total expectation of Lemma~\ref{lem:descent} and sum from $k=0$ to $T-1$:
\[
\E[f(x_T)]\le f(x_0)-\frac{\eta}{2}\sum_{k=0}^{T-1}\E\bigl[\|\nabla f(x_k)\|^{2}\bigr]
+2\eta L^{2}\sum_{k=0}^{T-1}\E[D_k].
\tag{7}
\]

\item \textbf{Evolution of the Distance Terms $D_k$.}  
Recall $D_k:=\frac1n\sum_{i=1}^{n}\|x_k-\phi_i^k\|^{2}$.  When index $i_k$ is sampled, only $\phi_{i_k}$ is updated to $x_k$; all other table entries stay unchanged.  Hence
\[
\begin{aligned}
\E[D_{k+1}\mid\mathcal{F}_k]
&=\frac1n\Bigl[\sum_{i\neq i_k}\|x_{k+1}-\phi_i^k\|^{2}+\|x_{k+1}-x_k\|^{2}\Bigr] \\[2pt]
&=\frac{n-1}{n}\E\bigl[\|x_{k+1}-\phi_i^k\|^{2}\bigr]_{i\neq i_k}
+\frac1n\|x_{k+1}-x_k\|^{2}.
\end{aligned}
\]
Taking expectation over $i_k$ and using symmetry of the $n-1$ unchanged indices gives
\[
\E[D_{k+1}]
\le \E[D_k]-\frac{1}{n}\E\bigl[\|x_k-\phi_{i_k}^k\|^{2}\bigr]
+\E\bigl[\|x_{k+1}-x_k\|^{2}\bigr].
\tag{8}
\]
But $\|x_{k+1}-x_k\|^{2}=\eta^{2}\|v_k\|^{2}$ and Lemma~\ref{lem:second} yields
\[
\E\bigl[\|x_{k+1}-x_k\|^{2}\bigr]
\le \eta^{2}\bigl(4L^{2}\E[D_k]+2\E\|\nabla f(x_k)\|^{2}\bigr).
\tag{9}
\]

\item \textbf{Bounding the Decrease of $D_k$.}  
Insert (9) into (8) and use $\E[\|x_k-\phi_{i_k}^k\|^{2}]=\E[D_k]$:
\[
\E[D_{k+1}]
\le \Bigl(1-\frac1n+4\eta^{2}L^{2}\Bigr)\E[D_k]
+\;2\eta^{2}\E\|\nabla f(x_k)\|^{2}.
\tag{10}
\]

\item \textbf{Summation of $D_k$ Recurrence.}  
Iterate (10) and sum the resulting inequality over $k=0,\dots,T-1$.  Because $1-\frac1n+4\eta^{2}L^{2}\le 1-\frac1{2n}$ for $\eta\le 1/(3L)$ (simple algebra), we obtain the geometric decay bound
\[
\sum_{k=0}^{T-1}\E[D_k]\;\le\; n\,\E[D_0]\;+\;2n\eta^{2}\sum_{k=0}^{T-1}\E\|\nabla f(x_k)\|^{2}.
\tag{11}
\]

\item \textbf{Plugging $D_k$ Bound into Descent Inequality.}  
Replace $\sum_{k=0}^{T-1}\E[D_k]$ in (7) using (11):
\[
\E[f(x_T)]\le f(x_0)-\frac{\eta}{2}\sum_{k=0}^{T-1}\E\|\nabla f(x_k)\|^{2}
+2\eta L^{2}\Bigl( n\,\E[D_0]+2n\eta^{2}\sum_{k=0}^{T-1}\E\|\nabla f(x_k)\|^{2}\Bigr).
\tag{12}
\]

\item \textbf{Collecting Gradient Norm Terms.}  
Move the gradient‑norm terms to the left-hand side:
\[
\Bigl(\frac{\eta}{2}-4\eta^{3}L^{2}n\Bigr)\sum_{k=0}^{T-1}\E\|\nabla f(x_k)\|^{2}
\le f(x_0)-\E[f(x_T)]+2\eta L^{2}n\,\E[D_0].
\tag{13}
\]
Because $\eta\le 1/(3L)$, the coefficient satisfies
\[
\frac{\eta}{2}-4\eta^{3}L^{2}n\;\ge\;\frac{\eta}{4},
\]
provided $n\ge1$ (the worst case $n=1$ still respects the inequality).  Hence
\[
\frac{\eta}{4}\sum_{k=0}^{T-1}\E\|\nabla f(x_k)\|^{2}
\le f(x_0)-\E[f(x_T)]+2\eta L^{2}n\,\E[D_0].
\tag{14}
\]

\item \textbf{Bounding the Table Initialization Term.}  
Recall $D_0=\frac1n\sum_{i=1}^{n}\|x_0-\phi_i^0\|^{2}$.  With the standard initialization $\phi_i^0=x_0$, we have $D_0=0$.  Even if the table is arbitrary, smoothness implies
\[
\|x_0-\phi_i^0\|\le \frac{1}{L}\|\nabla f_i(x_0)-\nabla f_i(\phi_i^0)\|
\]
and hence $D_0\le \frac{1}{L^{2}n}\sum_{i}\|\nabla f_i(x_0)-\nabla f_i(\phi_i^0)\|^{2}\le \frac{\sigma^{2}}{L^{2}}$, where $\sigma^{2}$ is defined in Assumption~\ref{as:bounded}.  For simplicity we adopt the common choice $\phi_i^0=x_0$ so that $D_0=0$ and the last term in (14) vanishes.

\item \textbf{Averaging Over Iterates.}  
Divide (14) by $T$ and define $\tilde x$ uniformly at random from $\{x_0,\dots,x_{T-1}\}$:
\[
\E\!\bigl[\|\nabla f(\tilde x)\|^{2}\bigr]
= \frac1T\sum_{k=0}^{T-1}\E\|\nabla f(x_k)\|^{2}
\le \frac{4\bigl(f(x_0)-\E[f(x_T)]\bigr)}{\eta T}.
\tag{15}
\]

\item \textbf{Introducing the Variance Term $\sigma^{2}$.}  
When $D_0\neq0$, the extra term $2\eta L^{2}n\,\E[D_0]$ in (14) becomes $2\eta L^{2}n\cdot\frac{\sigma^{2}}{L^{2}}=2\eta n\sigma^{2}$.  Since $n\ge1$, we may bound it by $2\eta\sigma^{2}$.  Adding this contribution to (15) yields
\[
\E\!\bigl[\|\nabla f(\tilde x)\|^{2}\bigr]
\le \frac{4\bigl(f(x_0)-f(x^{\star})\bigr)}{\eta T}+4L\eta\,\sigma^{2}.
\tag{16}
\]

\item \textbf{Final Form of the Bound.}  
Equation (16) is precisely the statement of Theorem~\ref{thm:main} (up to a harmless factor of $2$).  Renaming constants gives the compact expression (4).  ∎
\end{enumerate}

\section*{Rate Derivation (With Constants)}
From (4) set $\eta = \frac{1}{\sqrt{T}} \wedge \frac{1}{3L}$ (the minimum of the two).  
If $T\ge (3L)^{2}$, then $\eta = 1/(3L)$ and
\[
\E\|\nabla f(\tilde x)\|^{2}\le
\frac{6L\bigl(f(x_0)-f(x^{\star})\bigr)}{\sqrt{T}}
+ \frac{4}{3}\sigma^{2}.
\]
If $T<(3L)^{2}$ we are in the transient regime and the bound is dominated by the $1/(\eta T)$ term, which is still $O(1/\sqrt{T})$.  Hence the overall rate is
\[
\E\|\nabla f(\tilde x)\|^{2}=O\!\left(\frac{1}{\sqrt{T}}\right).
\]

\section*{Special Cases}
\begin{enumerate}
\item \textbf{Convex $L$‑smooth case.}  When each $f_i$ is convex, the same proof yields the stronger guarantee
\[
\E\bigl[f(\tilde x)-f(x^{\star})\bigr]\le O\!\left(\frac{1}{T}\right),
\]
matching the classical SAGA rate for convex optimisation.
\item \textbf{Full‑gradient limit.}  If $n=1$, the estimator $v_k$ reduces to the true gradient, and (4) collapses to the standard gradient‑descent bound $\E\|\nabla f(\tilde x)\|^{2}\le 2\bigl(f(x_0)-f(x^{\star})\bigr)/(\eta T)$.
\item \textbf{Large‑batch regime.}  Setting the batch size to $b>1$ (sampling $b$ indices per iteration) simply scales the variance term by $1/b$, yielding
\[
\E\|\nabla f(\tilde x)\|^{2}\le\frac{2\bigl(f(x_0)-f(x^{\star})\bigr)}{\eta T}
+\frac{4L\eta}{b}\sigma^{2}.
\]
\end{enumerate}

\end{document}